\chapter{Models and Data}\label{ch:b1:models-and-data}

On a Friday at the fish market of Brenmoor the price of cod was higher
than on Thursday, and more of it was sold. A reader who has drawn a
demand curve knows that a higher price should mean fewer sales. Either the
demand for fish slopes up, or the two days are not two points on one
demand curve. This chapter is about telling those two answers apart. It
sets out what a model is in this book --- a map from the variables it
takes as given to the variables it explains --- and the bookkeeping that
every later chapter uses: \omterm{def:b1:models-and-data:stock-flow}{stocks} and \omterm{def:b1:models-and-data:stock-flow}{flows}, growth rates, real and nominal
values, and the notation of the series. It ends with the oldest trap in
the reading of economic data: a cloud of prices and quantities is not a
curve.

\begin{openingchart}{The labour market as one price and one quantity: the yearly
change in total hours worked against the yearly change in labour compensation
per hour (deflated by the GDP deflator), United States, 1981--2023, in per cent.
Each dot is a year.}{Penn World Table 11.0; hours = persons employed $\times$
average annual hours; compensation per hour = labour share $\times$ real GDP at
constant national prices $\div$ hours}
\begin{tikzpicture}
\begin{axis}[omscatter, xlabel={Change in real compensation per hour (\%)},
  ylabel={Change in hours worked (\%)}, xmin=-3, xmax=8, ymin=-8, ymax=6,
  height=0.55\linewidth]
  \addplot table[x=dwage, y=dhours] {figdata/bachelor-1/01-models-and-data/out/labour_cloud.csv};
  \node[font=\footnotesize, anchor=west] at (axis cs:1.5,-5.9) {2009};
  \node[font=\footnotesize, anchor=east] at (axis cs:6.6,-6.7) {2020};
  \node[font=\footnotesize, anchor=west] at (axis cs:-0.3,4.4) {2021};
\end{axis}
\end{tikzpicture}
\end{openingchart}

The chart treats the American labour market as one market with one price,
the real pay of an hour of work, and one quantity, the hours worked. Its
43 dots do not lie on a curve. Their correlation is $-0.42$ and the line
that fits them best has a slope of $-0.53$: a year in which pay per hour
rose by one per cent more than usual was, on average, a year in which hours
grew by half a per cent less. % figdata: labour_cloud_stats
Read as a demand curve for labour, that is a gentle downward slope. Read as
a supply curve, it says that people work less when they are paid more. Both
readings are wrong in the same way: each dot is the meeting point of a
demand curve and a supply curve in one year, and both curves moved every
year. The 2020 dot makes the point sharply. Hours fell by 6.7 per cent and
pay per hour rose by 6.8 per cent, % figdata: labour_cloud
the largest rise of the period, in the year of the steepest fall in hours.
\Cref{sec:b1:models-and-data:identification} gives the model that reads
such a cloud, and what it takes to recover a curve from it.

\begin{recall}
A model keeps the few causes that matter for one question and states what
would prove it wrong; a fact is not a value judgement, and a correlation is
not a cause. A growth rate is a relative change, and a quantity growing at
a constant rate grows by equal ratios. A price index compares the cost of a
basket in two periods; a real value is a nominal value divided by a price
index. The logarithm turns products into sums: $\ln(xy) = \ln x + \ln y$.
\end{recall}

\section{Given and explained variables}\label{sec:b1:models-and-data:variables}

A model in this book has the frame the series has used since the school
volumes: what it assumes, what it predicts, where it was tested, and where
it breaks. With calculus at hand the first two fields can be made exact. A
model is a set of equations. Some of its variables are explained by it,
others are fed to it from outside, and a few numbers fix the shape of the
equations.

\begin{definition}[Endogenous and exogenous variables, parameters]\label{def:b1:models-and-data:variables}
\index{endogenous variable}\index{exogenous variable}\index{parameter}
A model is a system of equations $F(y, x; \theta) = 0$ that, for each value
of $x$ and $\theta$, has a solution $y = g(x; \theta)$. The variables $y$
that the model solves for are its \emph{endogenous variables}. The variables
$x$ that it takes as given, and that change from one observation to the
next, are its \emph{exogenous variables}. The numbers $\theta$ that fix the
shape of the equations and are the same for every observation are its
\emph{parameters}. The function $g$ is the model's prediction.
\end{definition}

The same variable can be endogenous in one model and exogenous in another.
In a model of the fish market, the weather at sea is exogenous: nothing in
the market changes it. The price of fish is endogenous. In a model of a
household's grocery budget the price of fish is exogenous, because one
household cannot move it. The choice is part of the model's assumptions,
and it is tested like any other: if the weather turned out to depend on
the market, the model would be wrong.

\begin{omfigure}
\begin{tikzpicture}[econ, every node/.style={font=\small}]
  \node[draw, thick, rounded corners, minimum width=38mm, minimum height=14mm,
        align=center] (m) at (0,0) {model\\[1pt] $F(y, x; \theta) = 0$};
  \node[align=center] (x) at (-4.6,0) {exogenous $x$\\(given)};
  \node[align=center] (y) at (4.6,0) {endogenous $y$\\(explained)};
  \node[align=center] (t) at (0,1.9) {parameters $\theta$\\(fixed)};
  \draw[->, thick] (x) -- (m);
  \draw[->, thick] (m) -- (y);
  \draw[->, thick] (t) -- (m);
\end{tikzpicture}

{\small A model as a map: the \omterm{def:b1:models-and-data:variables}{exogenous variables} and the \omterm{def:b1:models-and-data:variables}{parameters} go in,
the \omterm{def:b1:models-and-data:variables}{endogenous variables} come out, $y = g(x; \theta)$.}
\end{omfigure}

The market model is the first example. Demand and supply are linear, and
each is moved by a term the model does not explain:
\begin{equation}\label{eq:b1:models-and-data:market}
Q^d = a - b\,p + u^d, \qquad Q^s = c + d\,p + u^s, \qquad Q^d = Q^s = q,
\end{equation}
with $b > 0$ and $d > 0$. The \omterm{def:b1:models-and-data:variables}{endogenous variables} are $p$ and $q$. The
\emph{shocks} $u^d$ and $u^s$ are exogenous: they stand for everything that
moves demand or supply in a given period, such as the weather, a festival,
the price of feed or a change in tastes. The \omterm{def:b1:models-and-data:variables}{parameters} are $a$, $b$, $c$
and $d$.

\begin{definition}[Equilibrium of a model]\label{def:b1:models-and-data:equilibrium}
\index{equilibrium (of a model)}
An \emph{equilibrium} of a model is a value of its \omterm{def:b1:models-and-data:variables}{endogenous variables} at
which all of its equations hold at once, for given values of the \omterm{def:b1:models-and-data:variables}{exogenous
variables} and the \omterm{def:b1:models-and-data:variables}{parameters}. When the model describes a process that
repeats, $y_{t+1} = h(y_t; x, \theta)$, an equilibrium is a fixed point
$y^* = h(y^*; x, \theta)$: once there, the model stays there.
\end{definition}

\begin{proposition}[Equilibrium of the linear market]\label{prop:b1:models-and-data:linear}
\index{market model!linear}
If $b + d > 0$, the model \eqref{eq:b1:models-and-data:market} has exactly one
\omterm{def:b1:models-and-data:equilibrium}{equilibrium},
\[
p^* = \frac{a - c + u^d - u^s}{b + d}, \qquad
q^* = \frac{ad + bc + d\,u^d + b\,u^s}{b + d}.
\]
\end{proposition}

\begin{proof}
Setting $Q^d = Q^s$ gives $a - bp + u^d = c + dp + u^s$, a linear equation
in $p$ with the coefficient $b + d \neq 0$, so it has the single root
$p^*$. Then $q^* = c + d p^* + u^s
= \bigl[c(b + d) + d(a - c + u^d - u^s) + (b + d)u^s\bigr]/(b + d)$, which
simplifies to the expression stated.
\end{proof}

The proposition is the model's prediction, $g$ of
\cref{def:b1:models-and-data:variables}. Its derivatives are the
comparative statics of the market, met in the school volumes:
$\partial p^*/\partial u^d = 1/(b + d) > 0$ and
$\partial q^*/\partial u^d = d/(b + d) > 0$, so a rise in demand raises both
the price and the quantity, while $\partial p^*/\partial u^s < 0 <
\partial q^*/\partial u^s$. \Cref{ch:b1:comparative-statics} does this for
any model, linear or not.

\begin{example}[The fish market in numbers]\label{ex:b1:models-and-data:fish}
At Brenmoor, with $p$ in dalens per kilogram and $q$ in crates sold per
day, take $Q^d = 180 - 10p + u^d$ and $Q^s = 20 + 10p + u^s$ (invented
numbers). On a day with no shocks, $p^* = 160/20 = 8$ and
$q^* = (180 \times 10 + 10 \times 20)/20 = 100$. A storm that keeps 30
crates of supply in harbour, $u^s = -30$, raises the price to
$p^* = 190/20 = 9.5$ and lowers the quantity to
$q^* = (2000 - 300)/20 = 85$. A festival that brings 30 crates of extra
demand, $u^d = 30$, gives $p^* = 9.5$ and $q^* = 115$: a higher price and
more fish sold. That is Friday at the market.
\end{example}

\section{Stocks, flows, levels and rates}\label{sec:b1:models-and-data:stocks}

\begin{definition}[Stock, flow]\label{def:b1:models-and-data:stock-flow}
\index{stock}\index{flow}
A \emph{stock} is a quantity measured at a date: it has a unit with no time
in it (dalens, crates, machines, people). A \emph{flow} is a quantity
measured over a period: its unit is per unit of time (dalens per year,
crates per day). A stock changes only through flows: over a period, the
change in a stock equals its inflows minus its outflows.
\end{definition}

Wealth, the capital of a country, the money in circulation, public debt
and the number of people out of work are \omterm{def:b1:models-and-data:stock-flow}{stocks}. Income, investment, GDP,
the government deficit and the number of people who lose their jobs each
month are \omterm{def:b1:models-and-data:stock-flow}{flows}. The classic error is to compare the two: a debt of
\dal{236818} million is not ``larger'' than a GDP of \dal{393264} million in
the same sense as one GDP is larger than another, % (Brenland 2025, accounts.csv: debt, gdp; invented)
because the first is an amount and the second an amount per year. Their
ratio, 60 per cent of a year's GDP, is a number with a time unit in it.

\begin{omfigure}
\begin{tikzpicture}[econ]
  \fill[omDef!14] (0,0) rectangle (4,1.6);
  \draw[very thick] (0,2.6) -- (0,0) -- (4,0) -- (4,2.6);
  \draw[omDef, thick] (0,1.6) -- (4,1.6);
  \node at (2,0.8) {stock $K_t$};
  \draw[->, very thick, omThm] (-1.6,3.2) node[above, black] {inflow $I_t$}
        to[out=0, in=100] (0.9,1.9);
  \draw[->, very thick, omThm] (3.6,0.3) to[out=0, in=180] (5.6,-0.4)
        node[right, black] {outflow $\delta K_t$};
\end{tikzpicture}

{\small \omterm{def:b1:models-and-data:stock-flow}{Stock} and \omterm{def:b1:models-and-data:stock-flow}{flow}: the capital \omterm{def:b1:models-and-data:stock-flow}{stock} $K_t$ is the water in the tub;
investment $I_t$ runs in, depreciation $\delta K_t$ runs out, and
$K_{t+1} = (1 - \delta)K_t + I_t$.}
\end{omfigure}

The accumulation of capital is the example used most often in this book:
\begin{equation}\label{eq:b1:models-and-data:accumulation}
K_{t+1} = (1 - \delta)K_t + I_t ,
\end{equation}
where $\delta$ is the share of the capital \omterm{def:b1:models-and-data:stock-flow}{stock} that wears out in a year.
If investment stays at $I$ forever, the \omterm{def:b1:models-and-data:stock-flow}{stock} settles where inflow equals
outflow, $\delta K^* = I$, so $K^* = I/\delta$: an \omterm{def:b1:models-and-data:equilibrium}{equilibrium} in the sense
of \cref{def:b1:models-and-data:equilibrium}, a fixed point of
\eqref{eq:b1:models-and-data:accumulation}. With $I = 30$ and
$\delta = 0.05$, $K^* = 600$, twenty years of investment. The Solow model
of \cref{ch:b1:solow-model} is this equation with investment made
endogenous.

A growth rate turns a level into a rate. The school volumes measured it as
the relative change from one year to the next, $(x_{t+1} - x_t)/x_t$. With
calculus the cleaner object is the rate at an instant.

\begin{definition}[Continuous growth rate]\label{def:b1:models-and-data:growth-rate}
\index{continuous growth rate}\index{growth rate!continuous}
The \emph{continuous growth rate} of a positive differentiable quantity
$x(t)$ is
\[
g_x = \frac{\dot x}{x} = \frac{\dd \ln x}{\dd t}, \qquad \dot x = \frac{\dd x}{\dd t}.
\]
Between two dates $t_0 < t_1$ its average is
$\bar g_x = \bigl[\ln x(t_1) - \ln x(t_0)\bigr]/(t_1 - t_0)$, and a quantity
with a constant continuous growth rate $g$ is $x(t) = x(t_0)\,\eu^{g(t - t_0)}$.
\end{definition}

\begin{proposition}[Growth of a product and of a ratio]\label{prop:b1:models-and-data:product}
\index{growth rate!of a product}
For positive differentiable $x$ and $z$, and any constant $\alpha$,
\[
g_{xz} = g_x + g_z, \qquad g_{x/z} = g_x - g_z, \qquad g_{x^\alpha} = \alpha\, g_x .
\]
\end{proposition}

\begin{proof}
$\ln(xz) = \ln x + \ln z$, $\ln(x/z) = \ln x - \ln z$ and
$\ln x^\alpha = \alpha \ln x$; differentiate each with respect to $t$.
\end{proof}

So output per worker $Y/L$ grows at the growth rate of output minus that of
employment, and nominal GDP grows at the growth rate of real GDP plus that
of the price level. These rules are exact for continuous rates. For yearly
rates they are approximations, and the size of the error is worth knowing.

\begin{proposition}[Yearly and continuous rates]\label{prop:b1:models-and-data:log-approx}
\index{growth rate!yearly and continuous}
A quantity that grows by the yearly rate $g$, $x_{t+1} = (1 + g)x_t$, has
the \omterm{def:b1:models-and-data:growth-rate}{continuous growth rate} $\ln(1 + g)$, and for $|g| < 1$
\[
\ln(1 + g) = g - \frac{g^2}{2} + \frac{g^3}{3} - \dotsb ,
\qquad 0 \le g - \ln(1 + g) \le \frac{g^2}{2} \text{ for } g \ge 0 .
\]
\end{proposition}

\begin{proof}
$\ln x_{t+1} - \ln x_t = \ln(1 + g)$, so the average continuous rate over
the year is $\ln(1 + g)$. The series is the Taylor expansion of $\ln(1 + g)$
at $0$. For the bound, $h(g) = g - \ln(1 + g)$ has $h(0) = 0$ and
$h'(g) = g/(1 + g) \in [0, g]$ for $g \ge 0$; integrating from $0$ gives
$0 \le h(g) \le g^2/2$.
\end{proof}

At $g = 1\%$ the two rates differ by less than $0.005$ points
($\ln 1.01 = 0.00995$); at $5\%$ by $0.12$ points ($\ln 1.05 = 0.0488$); at
$40\%$ by more than six points ($\ln 1.40 = 0.336$). For the growth of a
rich economy the difference is noise. For an inflation of 40 per cent it is
not, and \cref{sec:b1:models-and-data:real} returns to it.

A logarithmic axis draws $\ln x$ against time while labelling the axis
with $x$ itself. A constant growth rate then draws a straight line, and the
slope of the line is the growth rate: the line climbs by the same distance
whenever $x$ is multiplied by the same factor.

\begin{omchart}{GDP per person, United Kingdom and Japan, 1870--2022, in 2011
dollars at purchasing-power parity, on a logarithmic axis. Each tick is ten
times the one below it, so a steeper line is a faster growth rate.}{Maddison
Project Database 2023; GDP per person}
\begin{tikzpicture}
\begin{semilogyaxis}[omtimesten, xlabel={Year}, ylabel={Dollars per person},
  ymin=1000, ymax=60000, xmin=1870, xmax=2022, legend pos=north west]
  \addplot table[x=year, y=GBR] {figdata/bachelor-1/01-models-and-data/out/gbr_jpn.csv};
  \addplot table[x=year, y=JPN] {figdata/bachelor-1/01-models-and-data/out/gbr_jpn.csv};
  \legend{United Kingdom, Japan}
\end{semilogyaxis}
\end{tikzpicture}
\end{omchart}

Read the slopes. From 1950 to 1973 Japan's GDP per person rose from
\usd{3062} to \usd{18226}, an average \omterm{def:b1:models-and-data:growth-rate}{continuous growth rate} of
$\ln(18226/3062)/23 = 7.76\%$ a year, % figdata: gbr_jpn_growth
the steepest stretch of either line. Over the same years the United
Kingdom grew at $2.39\%$. % figdata: gbr_jpn_growth
By 2022 the two levels were within half a per cent of each other, \usd{38406}
and \usd{38269}. % figdata: gbr_jpn
On an ordinary axis the first eighty years of the chart would be squeezed
against the bottom and the slopes could not be compared at all.

\section{Real and nominal}\label{sec:b1:models-and-data:real}

A value in money, the nominal value, mixes two things: a quantity and the
price it is valued at. National accounts separate them by valuing each
year's goods at the prices of one base year, which gives real GDP. The
ratio of the two is the GDP deflator, met in the school volumes:
\[
\text{deflator}_t = 100 \times \frac{\text{nominal GDP}_t}{\text{real GDP}_t},
\qquad
\text{real GDP}_t = \frac{\text{nominal GDP}_t}{\text{deflator}_t/100}.
\]

\begin{omchart}{GDP per person in the United States, 1960--2024, in current
dollars (nominal) and in dollars of 2020 (real), on a logarithmic axis. The
lines cross in 2020, the base year; the vertical gap between them is the
price level.}{World Bank, World Development Indicators; GDP per capita in
current and in constant local currency units}
\begin{tikzpicture}
\begin{semilogyaxis}[omtimesten, xlabel={Year}, ylabel={Dollars per person},
  ymin=1000, ymax=100000, xmin=1960, xmax=2024, legend pos=south east]
  \addplot table[x=year, y=nominal] {figdata/bachelor-1/01-models-and-data/out/us_gdppc.csv};
  \addplot table[x=year, y=real] {figdata/bachelor-1/01-models-and-data/out/us_gdppc.csv};
  \legend{Nominal, Real (2020 dollars)}
\end{semilogyaxis}
\end{tikzpicture}
\end{omchart}

Nominal GDP per person in the United States was \usd{3000} in 1960 and
\usd{86170} in 2024, 28.7 times as much. % figdata: us_gdppc, us_gdppc_growth
In dollars of 2020 it went from \usd{20565} to \usd{72395}, 3.52 times as
much. % figdata: us_gdppc, us_gdppc_growth
Over the 64 years the average \omterm{def:b1:models-and-data:growth-rate}{continuous growth rates} were $5.25\%$ a year
for the nominal series, $1.97\%$ for the real series and $3.28\%$ for the
deflator, and $5.25 = 1.97 + 3.28$, % figdata: us_gdppc_growth
as \cref{prop:b1:models-and-data:product} says it must: nominal GDP is the
product of real GDP and the deflator.

\begin{proposition}[Real growth]\label{prop:b1:models-and-data:real-growth}
\index{real growth}
If a nominal value grows by the yearly rate $g_N$ and the price index by
the yearly rate $\pi$, the real value grows by the yearly rate
\[
g_R = \frac{1 + g_N}{1 + \pi} - 1 = \frac{g_N - \pi}{1 + \pi},
\]
and the continuous rates add exactly: $\ln(1 + g_N) = \ln(1 + g_R) + \ln(1 + \pi)$.
The rule $g_R \approx g_N - \pi$ overstates $g_R$ by $\pi\, g_R$, the
product of the two rates.
\end{proposition}

\begin{proof}
The real value is $R_t = N_t/P_t$, so
$R_{t+1}/R_t = (N_{t+1}/N_t)/(P_{t+1}/P_t) = (1 + g_N)/(1 + \pi)$.
Taking logs gives the second statement. For the third, $1 + g_N =
(1 + g_R)(1 + \pi) = 1 + g_R + \pi + g_R\pi$, so
$g_N - \pi = g_R + \pi g_R$.
\end{proof}

\begin{measure}[Real values and the deflator]
A statistical office does not divide nominal GDP by a price index it chose
in advance. It values each year's quantities, product by product, at the
prices of a base year (or, in a chained index, at the prices of the year
before, linking the years together), and the deflator is what is left:
nominal divided by real. The base year is a choice. In the World Bank table
used for the chart, the American real series is in dollars of 2020, so
nominal and real are equal in 2020 and the deflator is 100 there; it was
14.59 in 1960 and 119.03 in 2024. % figdata: us_gdppc
A change of base year changes the level of every real number and leaves
\omterm{prop:b1:models-and-data:real-growth}{real growth} rates nearly unchanged; it is the ratios that carry
information. Real values of different countries are a second, separate
problem: converting them to one currency needs prices of the same goods in
both places (\cref{ch:b1:national-accounts-in-depth}).
\end{measure}

\section{The series notation}\label{sec:b1:models-and-data:notation}

The books of this series that use algebra use one notation, set out in
\cref{tab:b1:models-and-data:notation}. A symbol keeps its meaning across
chapters and across books. Where a letter is needed twice --- $M$ for
imports and for money, $u$ for utility, for unemployment and for a shock ---
the context separates them, and a subscript or superscript is added where
the two meet in one equation.

\begin{table}[ht]
\centering\small
\caption{The series notation.}\label{tab:b1:models-and-data:notation}
\begin{tabular}{@{}>{\raggedright\arraybackslash}p{0.30\linewidth}
                >{\raggedright\arraybackslash}p{0.48\linewidth}
                >{\raggedright\arraybackslash}p{0.14\linewidth}@{}}
\toprule
Symbol & Meaning & Used from \\
\midrule
$p$, $q$ & price and quantity of one good ($P$ is the price level) & this chapter \\
$Q^d$, $Q^s$ & quantity demanded and supplied ($D$, $S$ name the curves) & this chapter \\
$x_1$, $x_2$, $u(x_1, x_2)$ & the consumer's quantities and utility ($m$ income; $\lambda$ a Lagrange multiplier) & \cref{ch:b1:preferences-and-utility} \\
$\varepsilon$ & an elasticity ($\varepsilon_d$, $\varepsilon_s$); an error term is $u_i$ or $e_t$, never $\varepsilon$ & \cref{ch:b1:comparative-statics} \\
$Y$ & output ($y = Y/L$; $Y^*$ or $\bar Y$ potential output) & \cref{ch:b1:national-accounts-in-depth} \\
$C$, $I$, $G$, $X$, $M$, $NX$ & consumption, investment, government purchases, exports, imports, net exports ($T$ net taxes, $S$ saving) & \cref{ch:b1:national-accounts-in-depth} \\
$K$, $L$, $A$, $H$ & capital, labour, productivity, human capital ($k = K/L$; $\delta$ depreciation, $s$ saving rate, $n$ population growth, $g$ productivity growth) & this chapter \\
$\alpha$ & capital's share, in $Y = AK^\alpha L^{1-\alpha}$ & \cref{ch:b1:technology} \\
$w$, $W$ & real wage, nominal wage ($r$ is the rental rate in the firm's problem) & \cref{ch:b1:cost-minimisation} \\
$i$, $r$ & nominal and real interest rates ($r \approx i - \pi^e$; $i^*$ the foreign rate) & \cref{ch:b1:labour-supply-and-saving} \\
$\pi$, $\pi^e$ & inflation, expected inflation (profit is $\Pi$) & this chapter \\
$u$, $u^n$ & unemployment rate, natural rate & \cref{ch:b1:unemployment-flows} \\
$M$ & money supply ($M^s$ where it would clash with imports) & \cref{ch:b1:money-and-inflation} \\
$e$, $\varepsilon_r$ & nominal exchange rate (domestic currency per unit of foreign; a rise is a depreciation), real exchange rate & \cref{ch:b1:open-economy} \\
$\beta$ & the discount factor (Years 2--3); regression coefficients are $\beta_0$, $\beta_1$, with hats for estimates & \cref{ch:b1:simple-regression} \\
$t$, $\tau$ & time, a tax rate & this chapter \\
$\mathrm{MC}$, $\mathrm{AC}$, $\mathrm{AVC}$, $\mathrm{MR}$, $\mathrm{TR}$, $\mathrm{TC}$, $\mathrm{FC}$ & cost and revenue measures, set upright & \cref{ch:b1:cost-minimisation} \\
$\mathrm{CS}$, $\mathrm{PS}$, $\mathrm{DWL}$ & consumer surplus, producer surplus, deadweight loss & \cref{ch:b1:consumer-welfare} \\
\bottomrule
\end{tabular}
\end{table}

Two conventions in the table do work in later chapters. The exchange rate
is the price of foreign currency, so a rise in $e$ is a fall in the value of
the home currency (\cref{ch:b1:open-economy}). The relation
$r \approx i - \pi^e$, the real interest rate as the nominal rate less
expected inflation, is \cref{prop:b1:models-and-data:real-growth} applied
to a loan: a lender who is repaid $1 + i$ dalens when prices have risen by
$\pi$ gets back $(1 + i)/(1 + \pi)$ in goods (\cref{ch:b1:labour-supply-and-saving}).

\section{Observed prices and quantities}\label{sec:b1:models-and-data:identification}

Return to the opening chart. Each year, the market model
\eqref{eq:b1:models-and-data:market} produces one point $(p^*_t, q^*_t)$,
the meeting of that year's demand and supply curves. If only demand moved,
every point would sit on the one supply curve and the cloud would trace it.
If only supply moved, the cloud would trace demand. When both move, the
cloud traces neither.

\begin{omfigure}
\begin{tikzpicture}
\begin{axis}[omrecord, xlabel={Quantity $q$}, ylabel={Price $p$}, xmin=0, xmax=10,
  ymin=0, ymax=10, width=0.62\linewidth, height=0.55\linewidth, grid=none,
  clip=true]
  \addplot[demand] table[x=qd1, y=p] {figdata/bachelor-1/01-models-and-data/out/curves.csv};
  \addplot[demand] table[x=qd2, y=p] {figdata/bachelor-1/01-models-and-data/out/curves.csv};
  \addplot[demand] table[x=qd3, y=p] {figdata/bachelor-1/01-models-and-data/out/curves.csv};
  \addplot[supply] table[x=qs1, y=p] {figdata/bachelor-1/01-models-and-data/out/curves.csv};
  \addplot[supply] table[x=qs2, y=p] {figdata/bachelor-1/01-models-and-data/out/curves.csv};
  \addplot[supply] table[x=qs3, y=p] {figdata/bachelor-1/01-models-and-data/out/curves.csv};
  \addplot[only marks, mark=*, mark size=1.4pt, black] table[x=q, y=p]
    {figdata/bachelor-1/01-models-and-data/out/cloud.csv};
\end{axis}
\end{tikzpicture}

{\small Thirty draws of the market model with $a = 10$, $b = 1$, $c = 2$,
$d = 1$ and shocks of equal spread (simulated). Three of the thirty pairs of
curves are drawn, demand in blue, supply in red; the dots are the thirty
equilibria. The cloud shows almost no relation between price and quantity:
it is neither curve.}
\end{omfigure}

How weak is the relation? The line that fits a cloud of points best, in the sense of least
squares, has the slope
\[
\hat\beta = \frac{\sum_t (p_t - \bar p)(q_t - \bar q)}{\sum_t (p_t - \bar p)^2},
\]
where bars are averages over the observations
(\cref{ch:b1:simple-regression} derives it). Write
$\tilde u^d_t = u^d_t - \bar u^d$ for the deviation of a shock from its
average, and
\[
S_d = \sum_t (\tilde u^d_t)^2, \qquad S_s = \sum_t (\tilde u^s_t)^2, \qquad
S_{ds} = \sum_t \tilde u^d_t\, \tilde u^s_t .
\]

\begin{proposition}[The slope of the cloud]\label{prop:b1:models-and-data:cloud}
\index{statistical demand curve}
In the market model \eqref{eq:b1:models-and-data:market}, if the two shocks
are unrelated in the sense that $S_{ds} = 0$, the least-squares slope of $q$
on $p$ over the equilibria is
\[
\hat\beta = w\,d - (1 - w)\,b, \qquad w = \frac{S_d}{S_d + S_s} \in [0, 1].
\]
It is a weighted average of the supply slope $d$ and the demand slope $-b$,
the weights being the shares of the two shocks in the total variation.
\end{proposition}

\begin{proof}
By \cref{prop:b1:models-and-data:linear}, $p_t - \bar p = (\tilde u^d_t -
\tilde u^s_t)/(b + d)$ and $q_t - \bar q = (d\,\tilde u^d_t + b\,\tilde
u^s_t)/(b + d)$. Multiply and sum:
$\sum_t (p_t - \bar p)(q_t - \bar q) = (d S_d - b S_s + (b - d)S_{ds})/(b +
d)^2$, and $\sum_t (p_t - \bar p)^2 = (S_d + S_s - 2S_{ds})/(b + d)^2$. With
$S_{ds} = 0$ the ratio is $(d S_d - b S_s)/(S_d + S_s) = w d - (1 - w) b$.
\end{proof}

In the simulated figure $b = d = 1$ and the two shocks have nearly the same
spread, so $w$ is close to one half and the least-squares slope of $q$ on $p$
is close to zero: the
formula gives $-0.13$ and the fitted line $-0.18$. % figdata: cloud_stats
The gap is the term $S_{ds}$: in thirty draws the two shocks happened to
move together a little, and the full formula of
\cref{exo:b1:models-and-data:cross} gives $-0.18$ exactly. % figdata: cloud_stats
The proposition answers the question of the opening. The slope of a cloud
of prices and quantities tells you how large the demand shocks were
relative to the supply shocks. It does not tell you the slope of either
curve.

\begin{definition}[Identification problem]\label{def:b1:models-and-data:identification}
\index{identification problem}\index{identified parameter}
A \omterm{def:b1:models-and-data:variables}{parameter} of a model is \emph{identified} by a kind of data if no two
values of it, combined with the model's assumptions, could have produced
the same data. The \emph{identification problem} of a market is that
observed \omterm{def:b1:models-and-data:equilibrium}{equilibrium} prices and quantities alone identify neither the slope
of demand nor the slope of supply.
\end{definition}

What breaks the deadlock is a variable that moves one curve and not the
other.

\begin{definition}[Shifter]\label{def:b1:models-and-data:shifter}
\index{shifter}\index{demand shifter}\index{supply shifter}
A \emph{supply shifter} is an \omterm{def:b1:models-and-data:variables}{exogenous variable} that enters the supply
equation and not the demand equation. A \emph{demand shifter} is an
\omterm{def:b1:models-and-data:variables}{exogenous variable} that enters the demand equation and not the supply
equation.
\end{definition}

\begin{proposition}[Tracing a curve with a shifter]\label{prop:b1:models-and-data:trace}
\index{shifter!tracing a curve}
Split the observations of the market model into two groups by a \omterm{def:b1:models-and-data:shifter}{supply
shifter} $z \in \{0, 1\}$: supply is $Q^s = c + dp - \gamma z + u^s$. If the
demand shocks have the same average in both groups, $\bar u^d_1 = \bar
u^d_0$, and the average prices differ, then
\[
\frac{\bar q_1 - \bar q_0}{\bar p_1 - \bar p_0} = -b ,
\]
the slope of demand. Symmetrically, a \omterm{def:b1:models-and-data:shifter}{demand shifter} traces supply.
\end{proposition}

\begin{proof}
Each \omterm{def:b1:models-and-data:equilibrium}{equilibrium} lies on that period's demand curve: $q_t = a - b p_t +
u^d_t$, because $z$ does not enter demand. Average over each group: $\bar
q_g = a - b \bar p_g + \bar u^d_g$ for $g = 0, 1$. Subtract the two lines and
use $\bar u^d_1 = \bar u^d_0$: $\bar q_1 - \bar q_0 = -b(\bar p_1 - \bar
p_0)$. Divide by $\bar p_1 - \bar p_0 \neq 0$.
\end{proof}

Both assumptions carry weight. The \omterm{def:b1:models-and-data:shifter}{shifter} must stay out of the demand
equation, and the demand shocks must not differ between the two groups:
storms must not be more frequent on the days when buyers are fewer. Neither
can be checked from prices and quantities alone. They are claims about the
world, defended with knowledge of the market. The idea, an outside variable
that moves the market along one curve, is the instrument of the Year~3
volume, where it gets its general form.

\begin{model}{The market model with shifting curves}
\assumes Linear demand $Q^d = a - bp + u^d_t$ and supply $Q^s = c + dp +
u^s_t$; the price clears the market each period; the shocks to the two
curves are unrelated.
\predicts The least-squares slope of the cloud of equilibria is a weighted
average of the two slopes, weighted by the variation of the shocks; a
\omterm{def:b1:models-and-data:shifter}{supply shifter} traces demand, and a \omterm{def:b1:models-and-data:shifter}{demand shifter} traces supply.
\tested The yearly American labour-market cloud is consistent with it (a
correlation of $-0.42$ and a slope of $-0.53$, % figdata: labour_cloud_stats
neither curve), but a cloud alone cannot test it. The Fulton fish market
does: the fitted line has a slope of $-0.54$, while stormy days at sea, a
\omterm{def:b1:models-and-data:shifter}{supply shifter}, trace a demand slope of $-1.08$ (see the record box
below). % ledger: b1:graddy-ols, b1:graddy-iv
\breaks When the two shocks move together ($S_{ds} \neq 0$), as when one
event moves both curves; when the quantity is fixed before the price is
known, as for a crop planted months ahead; and when the price measured is
an average over a changing mix of goods or of workers.
\end{model}

The last breaking point applies to the opening chart. In April 2020 the
average hourly earnings of American private employees rose by \usd{1.34},
to \usd{30.01}, % ledger: b1:bls2020-ahe
and the statistical office that publishes them explained that the increase
largely reflected the job losses among lower-paid workers. % ledger: b1:bls2020-ahe
Pay per hour rose because the hours that disappeared were cheap ones, not
because the hour of any one worker cost more. The 2020 dot is not a point
of any one supply or demand curve; it is a change in the mix of workers.

\begin{record}[A fish market read twice]
From December 1991 to May 1992 one dealer's daily sales of whiting at the
Fulton fish market in New York were recorded over 111 market days (Graddy
1995). % ledger: b1:graddy-days
The least-squares line through the cloud of log quantities and log prices
has a slope of $-0.54$ (standard error 0.18): what an early analysis of
such fitted lines (Working 1927) called a ``\omterm{prop:b1:models-and-data:cloud}{statistical demand curve}.''
% ledger: b1:graddy-ols, b1:working1927
Bad weather at sea is a \omterm{def:b1:models-and-data:shifter}{supply shifter}: it keeps boats in harbour and does
not change how much whiting a fishmonger wants that morning. On stormy days
the log quantity was lower by 0.36 and the log price higher by 0.34, and
the ratio, about $-0.36/0.34 = -1.06$ with these rounded effects and
$-1.08$ (standard error 0.46) with the unrounded ones, is the slope of
demand that \cref{prop:b1:models-and-data:trace} recovers (Imbens 2014).
% ledger: b1:graddy-iv
The two numbers differ by a factor of two. Only the second has a reason to
be the demand curve, and its standard error shows how much the data still
leave open.
\end{record}

\section*{Exercises}\label{sec:b1:models-and-data:exercises}

\begin{exercise}[$\star$]\label{exo:b1:models-and-data:read-cloud}
In the opening chart, 2009 has hours $-5.9\%$ and pay per hour $+1.3\%$;
2021 has hours $+4.4\%$ and pay per hour $-0.4\%$. % figdata: labour_cloud
\begin{enumerate}
  \item If only one curve moved between 2008 and 2009, which one, and in which direction?
  \item The same question for 2020 to 2021.
  \item Give one reason, from the chapter, not to take the 2009 answer at face value.
\end{enumerate}
\end{exercise}

\begin{exercise}[$\star$]\label{exo:b1:models-and-data:per-worker}
Output grows at $3\%$ a year and employment at $1\%$.
\begin{enumerate}
  \item With continuous rates, at what rate does output per worker grow?
  \item With yearly rates, compute the exact rate $(1.03/1.01) - 1$ and compare.
\end{enumerate}
\end{exercise}

\begin{exercise}[$\star$]\label{exo:b1:models-and-data:brenland}
Brenland's nominal GDP was \dal{248438} million in 2015 and \dal{393264}
million in 2025; its GDP deflator, 2015 = 100, was 134.5 in 2025
(invented data, \code{accounts.csv}).
\begin{enumerate}
  \item Compute real GDP in 2025, in millions of 2015 dalens.
  \item Compute the average \omterm{def:b1:models-and-data:growth-rate}{continuous growth rates} of nominal GDP, real GDP and the deflator over 2015--2025.
  \item Check that the first is the sum of the other two.
\end{enumerate}
\end{exercise}

\begin{exercise}[$\star$]\label{exo:b1:models-and-data:stock-or-flow}
Say whether each is a \omterm{def:b1:models-and-data:stock-flow}{stock} or a \omterm{def:b1:models-and-data:stock-flow}{flow}, and give its unit: (a) the public
debt of Brenland; (b) the government deficit; (c) the number of people
unemployed; (d) the number of people who find a job in a month; (e) a
household's savings account balance; (f) a household's saving; (g) the
cod in the cold store of Brenmoor; (h) the catch landed at the harbour.
\end{exercise}

\begin{exercise}[$\star\star$]\label{exo:b1:models-and-data:high-inflation}
In a year in which nominal wages rise by $50\%$ and prices by $40\%$:
\begin{enumerate}
  \item compute the exact \omterm{prop:b1:models-and-data:real-growth}{real growth} of wages;
  \item compute the rule $g_N - \pi$ and the error it makes;
  \item compute the difference of the continuous rates, $\ln 1.5 - \ln 1.4$, and say which of the three numbers is the yearly \omterm{prop:b1:models-and-data:real-growth}{real growth}.
\end{enumerate}
\end{exercise}

\begin{exercise}[$\star\star$]\label{exo:b1:models-and-data:doubling}
\begin{enumerate}
  \item How long does a quantity growing at the constant continuous rate $2\%$ take to double?
  \item And one growing by $2\%$ every year?
  \item Compare both with the rule of 70, $70/2$, and explain where the 70 comes from.
\end{enumerate}
\end{exercise}

\begin{exercise}[$\star\star$]\label{exo:b1:models-and-data:slope-numbers}
In the market model take $b = 2$ and $d = 1$, and suppose the supply shocks
vary three times as much as the demand shocks: $S_s = 3S_d$, with $S_{ds} =
0$.
\begin{enumerate}
  \item Compute the weight $w$ and the least-squares slope of the cloud.
  \item Is the cloud closer to the demand curve or to the supply curve? Why?
  \item Predict how the slope changes if the supply shocks become smaller.
\end{enumerate}
\end{exercise}

\begin{exercise}[$\star\star$]\label{exo:b1:models-and-data:one-curve}
Predict the sign and the size of the cloud's slope, from
\cref{prop:b1:models-and-data:cloud}, when (a) only demand shifts and (b)
only supply shifts. For each case, say what the cloud traces, and what
observation in such a market would show the prediction wrong.
\end{exercise}

\begin{exercise}[$\star\star$]\label{exo:b1:models-and-data:log-chart}
On the logarithmic chart of the United Kingdom and Japan:
\begin{enumerate}
  \item The United Kingdom went from \usd{11061} in 1950 to \usd{19168} in 1973 and from \usd{19168} to \usd{31946} in 2000. % figdata: gbr_jpn_growth
    Compute the average growth rate of each period. Which is steeper on the chart?
  \item In 1913 the two countries stood at \usd{8212} and \usd{2431}. % figdata: gbr_jpn_growth
    By what factor was the United Kingdom ahead, and how does that factor show on a logarithmic axis?
  \item Why would an ordinary axis make the comparison of question 1 misleading?
\end{enumerate}
\end{exercise}

\begin{exercise}[$\star\star\star$]\label{exo:b1:models-and-data:cross}
Drop the assumption $S_{ds} = 0$ in \cref{prop:b1:models-and-data:cloud}.
\begin{enumerate}
  \item Derive the least-squares slope of the cloud in terms of $b$, $d$, $S_d$, $S_s$ and $S_{ds}$.
  \item Take $b = 2$, $d = 1$, $S_d = 1$, $S_s = 3$, $S_{ds} = 1$ and compute it.
  \item Name an event that would move both curves in the same direction on the same day, and say how it makes $S_{ds}$ positive.
\end{enumerate}
\end{exercise}

\begin{exercise}[$\star\star\star$]\label{exo:b1:models-and-data:bad-shifter}
At Brenmoor, $Q^d = 180 - 10p + u^d$ and $Q^s = 20 + 10p + u^s$. A storm
removes 30 crates of supply, and keeps buyers at home too, removing 10
crates of demand.
\begin{enumerate}
  \item Compute the change in price and in quantity that the storm causes.
  \item Compute the ratio of the two changes, as \cref{prop:b1:models-and-data:trace} would, and compare it with the demand slope.
  \item Which assumption of the proposition failed? What data, besides prices and quantities, would reveal the failure?
\end{enumerate}
\end{exercise}

\begin{exercise}[$\star\star\star$]\label{exo:b1:models-and-data:falsify}
Someone claims that the American labour market of the opening chart was
moved only by shifts in the demand for labour, so that the cloud traces a
fixed supply curve.
\begin{enumerate}
  \item What would the cloud look like if the claim were true?
  \item Name the feature of the chart that contradicts the claim, and the
    one dot that contradicts it most clearly.
  \item Now take the claim that only supply moved. What would the cloud
    look like, and what in the chart contradicts it?
\end{enumerate}
\end{exercise}

\section{Problem: The Fish Market of Brenmoor}\label{sec:b1:models-and-data:problem}

\begin{problem}[{Weekend problem --- forty days at the fish market}]\label{pb:b1:models-and-data:fish}
The fish market of Brenmoor (invented) sells cod by the crate, at a price
in dalens per kilogram; a crate holds 10~kilograms. Over forty days of the
2025 season the market kept a daily record of the price and the crates
sold. On ten days a storm kept part of the fleet in harbour. On four days,
a festival weekend in two parts, the town filled with visitors. The chart
shows the forty days and the table the sums over each kind of day.

\begin{omchart}{Forty days at the fish market of Brenmoor: price and crates
sold, by kind of day (invented data).}{the chapter's seeded script; no real data}
\begin{tikzpicture}
\begin{axis}[omscatter, xlabel={Price (dalens per kilogram)},
  ylabel={Crates sold}, xmin=7, xmax=10, ymin=75, ymax=120,
  width=0.7\linewidth, height=0.48\linewidth, legend pos=south west]
  \addplot table[x=price, y=crates] {figdata/bachelor-1/01-models-and-data/out/fish_calm.csv};
  \addplot table[x=price, y=crates] {figdata/bachelor-1/01-models-and-data/out/fish_storm.csv};
  \addplot table[x=price, y=crates] {figdata/bachelor-1/01-models-and-data/out/fish_fest.csv};
  \legend{Calm ordinary days, Storm days, Festival days}
\end{axis}
\end{tikzpicture}
\end{omchart}

\begin{center}\small
\pgfplotstabletypeset[col sep=comma, string type,
  columns={kind,days,sum_price,sum_crates},
  columns/kind/.style={column name={Kind of day}, column type={l}},
  columns/days/.style={column name={Days}, column type={r}},
  columns/sum_price/.style={column name={Sum of prices}, column type={r}},
  columns/sum_crates/.style={column name={Sum of crates}, column type={r}},
  every head row/.style={before row=\toprule, after row=\midrule},
  every last row/.style={after row=\bottomrule}]
  {figdata/bachelor-1/01-models-and-data/out/fish_groups.csv}
\end{center}

Keep prices to four decimals and crates to two when you compute averages.

\medskip\noindent\textbf{Part I --- The cold store.}
Unsold cod goes to the market's cold store, where $2\%$ of the \omterm{def:b1:models-and-data:stock-flow}{stock} spoils
each day. Let $S_t$ be the crates in store at the start of day $t$, $C_t$
the crates landed on day $t$ and $Q_t$ the crates sold. The cold store
holds 50 crates at the start of day 1. The boats land 110, 95 and 120 crates
on days 1 to 3, and the market sells 99, 99 and 80.
\begin{enumerate}
  \item Which of $S_t$, $C_t$ and $Q_t$ are \omterm{def:b1:models-and-data:stock-flow}{stocks} and which are \omterm{def:b1:models-and-data:stock-flow}{flows}? Give the units.
  \item Write the equation that takes $S_t$ to $S_{t+1}$.
  \item Compute $S_2$, $S_3$ and $S_4$, to one decimal.
  \item If, every day, two crates more are landed than are sold, at what level does the \omterm{def:b1:models-and-data:stock-flow}{stock} settle?
  \item Show that this level is an \omterm{def:b1:models-and-data:equilibrium}{equilibrium} in the sense of \cref{def:b1:models-and-data:equilibrium}.
  \item Starting from $S_1 = 50$, is the \omterm{def:b1:models-and-data:stock-flow}{stock} rising or falling towards that level? Justify.
\end{enumerate}

\medskip\noindent\textbf{Part II --- Real and nominal.}
In the same forty days of 2024 the average price was \dal{7.90} per
kilogram (invented). Over the forty days of 2025 it was \dal{8.50}, to two
decimals. Brenland's consumer price index was 140.0 in 2024 and 143.5 in
2025 (\code{accounts.csv}).
\begin{enumerate}[resume]
  \item Compute the nominal rise in the price of cod, in per cent, to one decimal.
  \item Compute the inflation rate of the consumer price index.
  \item Compute the real rise in the price of cod, exactly, and by the approximate rule.
  \item Express the 2025 price in 2024 dalens.
  \item A fisher says ``cod is 7.6~per cent dearer than last year.'' Is that true in real terms? Say what the right statement is.
\end{enumerate}

\medskip\noindent\textbf{Part III --- The cloud.}
\begin{enumerate}[resume]
  \item Compute the average price and the average crates on calm ordinary days, on storm days and on festival days.
  \item The least-squares slope of crates on price over all forty days is $-5.32$. % figdata: fish_stats
    Using \cref{prop:b1:models-and-data:cloud}, explain why this is not the slope of the demand for cod.
  \item Compare the festival days with the calm days. Which curve did the festival move, and in which direction?
  \item Compute the ratio of differences $(\bar q_{\text{fest}} - \bar q_{\text{calm}})/(\bar p_{\text{fest}} - \bar p_{\text{calm}})$. Which curve's slope does it estimate, and under which assumption?
  \item Compute the market's revenue, in dalens, on an average calm day and on an average storm day.
  \item Explain why the two revenues are close, using what you will find in Part IV.
\end{enumerate}

\medskip\noindent\textbf{Part IV --- Using the storms.}
\begin{enumerate}[resume]
  \item Explain why a storm is a \omterm{def:b1:models-and-data:shifter}{supply shifter} and not a \omterm{def:b1:models-and-data:shifter}{demand shifter} in this market. What would make it both?
  \item Compute the ratio of differences between storm days and calm days. Why are the festival days left out?
  \item Over the forty days the average price is \dal{8.4998} per kilogram and the average sale 98.25 crates. % figdata: fish_stats
    Compute the price elasticity of demand at these averages implied by the slope of question 19.
  \item Is demand elastic or inelastic at the averages? What does a storm then do to the fishers' total revenue?
  \item The script that made the forty days used a demand slope of exactly $-10$ crates per dalen. Compute the elasticity at the averages with that slope, and give one reason why question 19 did not recover $-10$ exactly.
  \item The storm days are the 3rd, 7th, 8th, 14th, 19th, 22nd, 23rd, 29th, 33rd and 38th. Name one way the assumption of \cref{prop:b1:models-and-data:trace} could still fail with such dates.
  \item State the chapter's result for this market in one sentence: the demand elasticity of cod at the mean, recovered from storm days, and what the plain cloud would have said instead.
\end{enumerate}
\end{problem}
